%% COULEUR
\usepackage{xcolor}

% 1. Choix de la matière
% maths, agregmusique, concoursBanqueFrance
\newcommand{\doctype}{DGFIP} % changer ici selon besoin

% 2. Définir une macro qui charge les couleurs selon la valeur de \doctype
\makeatletter
\newcommand{\loadcharte}{%
  \@ifundefined{doctype@colors@\doctype}{%
    % si type inconnu, couleur par défaut
    \definecolor{primarycolor}{HTML}{DA291C} %% Rouge
    \definecolor{secondarycolor}{HTML}{005EB8} %% Bleu
    \definecolor{background1}{HTML}{ffeeee} %% Rose
    \definecolor{background2}{HTML}{eaecee} %% Gris
  }{%
    % sinon appel à la macro spécifique
    \csname doctype@colors@\doctype\endcsname
  }%
}

% 3. Définir les macros pour chaque charte graphique

\newcommand{\doctype@colors@maths}{%
    \definecolor{primarycolor}{HTML}{C63C15}
    \definecolor{secondarycolor}{HTML}{005EB8}
    \definecolor{background1}{HTML}{FCECE7}
    \definecolor{background2}{HTML}{E8F1FA}
}

\newcommand{\doctype@colors@agregmusique}{%
    \definecolor{primarycolor}{HTML}{005EB8}
    \definecolor{secondarycolor}{HTML}{9754B8}
    \definecolor{background1}{HTML}{E8F1FA}
    \definecolor{background2}{HTML}{FAF3FD}
}

\newcommand{\doctype@colors@concoursBanqueFrance}{%
    \definecolor{primarycolor}{HTML}{AE2573}
    \definecolor{secondarycolor}{HTML}{005F61}
    \definecolor{background1}{HTML}{F7E7FA}
    \definecolor{background2}{HTML}{E5EFEF}
}

\newcommand{\doctype@colors@DGFIP}{%
    \definecolor{primarycolor}{HTML}{007681}
    \definecolor{secondarycolor}{HTML}{AE2573}
    \definecolor{background1}{HTML}{E5F1F2}
    \definecolor{background2}{HTML}{F7E8F1}
}

\loadcharte

\makeatother
